\documentclass[10pt,a4paper]{amsart}
\usepackage[margin=19mm]{geometry}
\usepackage{amsmath,amssymb,mathtools}
\usepackage[scr=rsfs,cal=euler]{mathalfa}
\usepackage{xfrac}
\usepackage[unicode,hidelinks,bookmarks=false]{hyperref}
\newcommand{\ie}{i.e.\ }
\newcommand{\cf}{cf.\ }
\newcommand{\resp}{respectively\ }
\newcommand{\Subsection}{\subsection}
\providecommand{\nobreakdash}{\nobreak\hbox{-}}
\setlength{\emergencystretch}{2em}
\multlinegap0pt
\usepackage{amssymb}
\usepackage[shortlabels]{enumitem}
\usepackage{float}
\usepackage{csquotes}
\usepackage{tikz}
\newtheorem{thm}{Theorem}
\def \cN {\mathcal{N}}
\DeclareMathOperator{\out}{out}
\def \sumt {\sum\limits}
\title{A series of real networks invariants}
\author{Mikhail Tuzhilin}
\date{Source: arXiv:2601.02052v2 (CC BY 4.0)}
\begin{document}
\maketitle

\noindent\emph{Source and licence.} A series of real networks invariants. Version arXiv:2601.02052v2 (CC BY 4.0), distributed by arXiv under the Creative Commons Attribution 4.0 International licence, http://creativecommons.org/licenses/by/4.0/, as recorded in the arXiv licence metadata for this exact version and shown on the abstract page https://arxiv.org/abs/2601.02052v2. Full licence notice: \texttt{licenses/arxiv-2601.02052-CC-BY-4.0.txt}.\par\smallskip

\noindent\textit{Editorial scope.} Counted unit: the article's thm thm1 (source0.tex lines 141--147). The article's complete original proof of it is delivered with it (source0.tex lines 148--150, one \texttt{proof} environment of 3 lines). No mathematics is written, repaired or re-derived: the statement and the proof are the article's own lines, in the article's own notation, and no retained line is substituted or re-flowed.  No further source-local context is needed: every cross-reference of every retained line resolves inside the two delivered spans.  Nothing was omitted from any retained span, so the package carries no omission record.  No retained line contains a citation, so the wrapper carries no bibliography and no bibliography entry.  No retained line contains a figure environment, an image-inclusion command or drawing code, so no image is delivered and the package declares no asset dependency.  Editorial formatting only, in addition: an AMS \texttt{amsart} preamble that restates the article's own declarations and macros used by the retained lines, namely the environments \texttt{thm} on separate counters, together with the standard packages those lines need.  The retained environments print in this document as Theorem 1 in order of appearance, whereas the article prints the same items with the numbering of its own full document.  The complete native source of the article as distributed in the arXiv e-print for this exact version is bundled unchanged as \texttt{sources/192132/source0.tex}.
\begin{thm}\label{thm1}
    For any graph $j$-neighborhood centrality
    $$
        \xi^j_i = \frac {\big|\out\big(\cN^j(i)\big)\Big|} {\big|\cN^j(i)\big|},
    $$
    where $\out(H) = E(H, V(G)\setminus H)$ the set of outer connections from the set $H$ to other vertices.
\end{thm}

\begin{proof}
    It is very simple fact from the formula for Laplacian matrix bilinear form: $\big(L x, x\big) = \sumt_{k,l\in V(G), k\sim l}(x^k-x^l)^2$ and thus $(x_k-x_l)^2 = 1$ only if $x_k = 1$ (and thus $k\in \cN^j(i)$) and $x_l = 0$ (and thus $l\notin \cN^j(i), l\sim k$) and vice versa.
\end{proof}

\end{document}
